\section{引言：从 Flow Matching 1 到 Flow Matching 2}

本讲是 KAIST CS492D（扩散模型与流模型）的第 10 讲，也是 Flow Matching 专题的第二部分。讲者 Minhyuk Sung 首先回顾了上一讲中引入的 Flow Matching 基本框架——流映射（Flow Map）、向量场（Vector Field）和概率路径（Probability Path）三大核心概念及其相互关系，然后在此基础上深入探讨以下关键问题：

\begin{enumerate}
    \item 扩散模型与流模型的\textbf{数学等价关系}：在何种条件下两者等价，如何互相转换？
    \item \textbf{速度场与分数函数的对应关系}：velocity 与 score function 之间的精确数学表达式是什么？
    \item \textbf{Rectified Flow}：如何通过直线轨迹简化生成过程，又如何通过 Reflow 微调使轨迹更直？
    \item \textbf{Optimal Transport} 视角：随机配对 vs.\ 最优配对如何影响轨迹弯曲程度？
    \item \textbf{实际应用}：Stable Diffusion 3（SD3）和 Flux 模型如何运用 Flow Matching 技术？
\end{enumerate}

\begin{knowledgebox}{前置知识回顾}
理解本讲需要以下来自上一讲（Lecture 9: Flow Matching 1）的前置概念：
\begin{itemize}
    \item \textbf{流映射（Flow Map）} $\psi_t$：从 $t=0$ 到 $t$ 的确定性映射，$\psi_0(\mathbf{x}) = \mathbf{x}$
    \item \textbf{向量场（Vector Field）} $u_t(\mathbf{x})$：流映射的时间导数，$\frac{d\psi_t}{dt} = u_t(\psi_t(\mathbf{x}))$
    \item \textbf{概率路径（Probability Path）} $p_t$：随时间 $t$ 变化的概率分布，满足 $p_0 = p_{\text{ref}}$, $p_1 = p_{\text{data}}$
    \item \textbf{Fokker--Planck 方程}：向量场与概率路径之间的联系方程
    \item \textbf{条件流匹配（Conditional Flow Matching）}：将边际向量场分解为对条件向量场的期望
\end{itemize}
\end{knowledgebox}

\subsection{符号约定与时间域的变化}

在 Flow Matching 的语境中，时间域与扩散模型有所不同：

\begin{itemize}
    \item \textbf{扩散模型}：时间从 $T$（噪声端）到 $0$（数据端），生成是``逆过程''
    \item \textbf{流模型}：时间从 $0$（参考分布）到 $1$（数据分布），生成是``正向过程''
\end{itemize}

具体地：
\begin{itemize}
    \item $\mathbf{x}_0 \sim p_0 = \mathcal{N}(\mathbf{0}, \mathbf{I})$：来自参考分布（标准高斯）的样本
    \item $\mathbf{x}_1 \sim p_1 = p_{\text{data}}$：来自数据分布的样本
\end{itemize}

\begin{warningbox}{容易混淆的符号}
注意在 Flow Matching 中 $\mathbf{x}_0$ 是\textbf{噪声}而 $\mathbf{x}_1$ 是\textbf{数据}，这与 DDPM 中 $\mathbf{x}_0$ 是数据、$\mathbf{x}_T$ 是噪声的惯例恰好相反。此外，Lecture 9 中的 $\mu_t, \sigma_t$ 函数与变分扩散模型中的 $\alpha_t, \sigma_t$ 本质相同，只是由于时间域翻转，位置互换了。
\end{warningbox}

\subsection{条件流映射的一般形式}

上一讲中引入了条件流映射（Conditional Flow Map）的线性高斯形式：
$$\psi_t(\mathbf{x}_0 \mid \mathbf{x}_1) = \sigma_t \cdot \mathbf{x}_0 + \alpha_t \cdot \mathbf{x}_1$$

其中：
\begin{itemize}
    \item $\sigma_t, \alpha_t$ 是关于时间 $t$ 的标量函数
    \item 边界条件：$\sigma_0 = 1, \alpha_0 = 0$（$t=0$ 时得到纯噪声 $\mathbf{x}_0$）；$\sigma_1 = 0, \alpha_1 = 1$（$t=1$ 时得到数据 $\mathbf{x}_1$）
\end{itemize}

对应的条件概率路径为高斯分布：
$$p_t(\mathbf{x} \mid \mathbf{x}_1) = \mathcal{N}(\mathbf{x}; \alpha_t \mathbf{x}_1, \sigma_t^2 \mathbf{I})$$

对条件流映射取时间导数，得到条件向量场：
$$u_t(\mathbf{x} \mid \mathbf{x}_1) = \dot{\sigma}_t \cdot \mathbf{x}_0 + \dot{\alpha}_t \cdot \mathbf{x}_1$$

\subsection{本章小结}

本节回顾了 Flow Matching 的基本设定，包括时间域约定、条件流映射的定义和条件向量场的推导。关键要点在于：流模型的时间从 $0$（噪声）到 $1$（数据），训练目标是学习一个神经网络来预测向量场（速度），这在概念上类似于扩散模型中预测噪声。

